\documentclass[a4paper,11pt]{article}

\usepackage{coworkingview}
\usepackage{hyperref}
\hypersetup{
  hidelinks,
  pdftitle={coworkingview},
  pdfauthor={CoworkingView},
}

\newcommand{\cmd}[1]{\texttt{\detokenize{#1}}}
\newcommand{\meta}[1]{$\langle$\textit{#1}$\rangle$}

\title{coworkingview}
\author{CoworkingView}
\date{Version 0.1.0}

\emergencystretch=2em

\begin{document}
\maketitle

\section{Purpose}

\cmd{coworkingview} prints median prices from a snapshot of the
CoworkingView Price Index.
The snapshot is shipped in
\begin{center}
\cmd{coworkingview-rates.csv}.
\end{center}
The package reads that file from the TeX input path.
It does not open a network connection and it does not run shell commands.

The code and this document are distributed under the LaTeX Project
Public License, version 1.3c or later.
The rates file is not part of that work.
It is Price Index data under CC BY 4.0
(\url{https://creativecommons.org/licenses/by/4.0/}).
Keep \cmd{\cwvattribution} with any figure taken from the snapshot.

\section{Snapshot}

The file records three cities that have a price row:
Abu Dhabi, Berlin and Madrid.
The unpublished fixture city \cmd{example-town} is in the project's
cities table and has no prices, so it is not in this file.
\cmd{\cwvrate} and \cmd{\cwvratestable} raise a package error for it,
as they do for any other unknown city.

Each row is copied from the open rates table.
The median, the percentiles and the sample size are the recorded
figures, not values computed by the package.
The snapshot date is the \cmd{as_of} value on those rows:

\begin{quote}
\cwvsnapshotdate
\end{quote}

The column \cmd{low_confidence} is kept in the file and is not printed
by the commands.
Empty \cmd{p25} and \cmd{p75} cells mean that percentile was not published.
Day passes and meeting rooms in this snapshot have no percentiles.

The recorded response credits the figures as follows.
Aggregated figures only; no endorsement by CoworkingView or any
operator is implied.

\begin{quote}
Source: CoworkingView Price Index
(\url{https://coworkingview.com/en/tools/price-index}),
as of October 2026.
Licensed under CC BY 4.0
(\url{https://creativecommons.org/licenses/by/4.0/}).
Aggregated figures only; no endorsement by CoworkingView or any
operator is implied.
\end{quote}

\section{Commands}

The package defines five commands.
An unknown city is an error whose message lists the cities in the
snapshot.
An unknown product is an error whose message lists the products
recorded for that city.

\subsection{\cmd{\cwvrate}}

\texttt{\string\cwvrate\{\meta{city}\}} prints the median, the currency
code and the unit for \cmd{flex_desk}.
\texttt{\string\cwvrate[product=\meta{product}]\{\meta{city}\}} selects
another product.
The text is the median as stored, a space, the currency code, a slash
and the unit, for example \cwvrate{berlin}.

\meta{city} is a slug or a display name.
Comparison ignores case.
Spaces, hyphens and underscores are the same separator, so
\cmd{Abu Dhabi}, \cmd{abu-dhabi} and \cmd{abu_dhabi} are one city.

\meta{product} is \cmd{flex_desk}, \cmd{fix_desk} or
\cmd{private_office}, or \cmd{day_pass} or \cmd{meeting_room}.
The same separator rule applies, so \cmd{private office} matches
\cmd{private_office}.
\cmd{flex_desk}, \cmd{fix_desk} and \cmd{private_office} are monthly.
\cmd{day_pass} is daily.
\cmd{meeting_room} is hourly.
The unit in the output is the unit from the row.

\begin{quote}
\small
\begin{tabular}{@{}ll@{}}
\texttt{\string\cwvrate\{berlin\}} & \cwvrate{berlin} \\
\texttt{\string\cwvrate[product=day\_pass]\{Berlin\}} &
  \cwvrate[product=day_pass]{Berlin} \\
\texttt{\string\cwvrate[product=private office]\{madrid\}} &
  \cwvrate[product=private office]{madrid} \\
\texttt{\string\cwvrate\{Abu Dhabi\}} & \cwvrate{Abu Dhabi} \\
\end{tabular}
\end{quote}

\subsection{\cmd{\cwvratestable}}

\texttt{\string\cwvratestable\{\meta{city}\}} prints a \cmd{tabular} with four
columns: product, median, the interval from the 25th to the 75th
percentile, and the sample size~$n$.
Products are ordered \cmd{flex_desk}, \cmd{fix_desk},
\cmd{private_office}, \cmd{day_pass}, \cmd{meeting_room}.
A product that is absent for that city is omitted.
A missing percentile is an em dash.
The percentiles use the currency shown in the median column.

Berlin:

\medskip
\cwvratestable{berlin}

\medskip
Madrid:

\medskip
\cwvratestable{Madrid}

\medskip
Abu Dhabi:

\medskip
\cwvratestable{abu-dhabi}

\subsection{Cities, date and credit}

\cmd{\cwvcities} prints the display names in the order of the file:
\cwvcities.

\cmd{\cwvsnapshotdate} prints \cwvsnapshotdate.

\cmd{\cwvattribution} prints the credit used next to a figure.
This document loads \cmd{hyperref}, so the price-index name is a link
and the visible text is
``Data: CoworkingView Price Index, CC BY 4.0''.
Without \cmd{hyperref}, the package does not load it.
The same words are printed and the URL is added with \cmd{\url}.

\begin{quote}
\cwvattribution
\end{quote}

\section{Files}

\begin{center}
\begin{tabular}{@{}ll@{}}
\cmd{coworkingview.sty} & package \\
\cmd{coworkingview-rates.csv} & snapshot, CC BY 4.0 \\
\cmd{coworkingview-doc.tex} & this document \\
\cmd{coworkingview-doc.pdf} & typeset document \\
\cmd{README.md} & short description and licence \\
\cmd{LICENSE} & LPPL 1.3c, and the data licence \\
\end{tabular}
\end{center}

\noindent Homepage\par
\url{https://coworkingview.com/en/tools/price-index}\par
\noindent Repository\par
\url{https://github.com/coworkingview/flex-office-data}\par
\noindent Bug reports\par
\url{https://github.com/coworkingview/flex-office-data/issues}

\end{document}
